\documentclass{nfr}

% This LaTeX template is based on the official Word template found here: https://www.forskningsradet.no/siteassets/utlysninger/vedlegg-utlysninger/template-for-project-description-researcher-project-for-experienced-scientists-and-early-career-scientists.docx
% See the README.md file for more information

% Here's an example bibliography using the numeric (Nature) style and two
% columns for the reference list.

% \usepackage{multicol}
% \renewcommand{\bibpreamble}{\begin{multicols*}{2}}
% \renewcommand{\bibpostamble}{\end{multicols*}}

% \usepackage{biblatex}
% \addbibresource{references.bib}

% \citestyle{nature}
% \citestyle{}

% \bibliographystyle{bst/veryshortnum}
\bibliographystyle{apalike}

\title{Research Plan}
\author{Zhongli Wang}

\begin{document}

\begin{center}
{\Large \textbf{Research Plan}}\\[0.6cm]

{\large \textbf{AI, Networks, and the Diffusion of Digital Technologies in Nordic Societies}}\\[0.4cm]
\end{center}

\noindent \textbf{Applicant:} Zhongli Wang\\
\textbf{Intended Host Department:} Economics / School of Management, University of Turku\\
\textbf{Secondment Partner:} Lucidminds AI (Netherlands)\\[0.6cm]

\section{Overview and Motivation}

Nordic societies are among the most digitally advanced in the world, yet the diffusion of emerging AI technologies — particularly household-level adoption and firm-level integration of large language models (LLMs) and automation tools — remains uneven. This heterogeneity is visible across firm size, sector, and region as well as across households by age, education, and digital skills, as documented in OECD digital readiness indicators and national statistics. Differences in adoption are shaped not only by economic incentives but also by the structure of social networks, organizational norms, trust in digital infrastructures, and the visibility of relevant use cases. As Finland implements ambitious national AI strategies and invests heavily in digital public services, understanding \emph{how new technologies spread through households, firms, and communities} is more crucial than ever.

Recent indicators suggest that gaps are economically meaningful: despite near-universal connectivity, measured AI uptake and advanced digital tool adoption remain substantially lower among SMEs and in lagging regions relative to frontier firms and sectors (e.g., OECD indicators; Statistics Finland/Tilastokeskus surveys). For a 2026 submission, the project will explicitly align measurement choices and policy scenarios with the most recent updates to Finland’s national AI and digitalization strategy (post-2025), ensuring timeliness and policy relevance.

This project develops a comprehensive research program to study the \textbf{diffusion of AI and digital technologies through social and economic networks in Nordic societies}, focusing on Finland as a core empirical setting. Finland provides a particularly informative case because it combines high-quality administrative and survey data infrastructure with high-trust institutions and advanced digital public services, while sharing core features with other Nordic economies (welfare-state institutions, high connectivity, and coordinated policy environments). The project integrates:

\begin{itemize}
\item \textbf{Micro-level behavior} (household AI use, firm digitalization choices)
\item \textbf{Network diffusion mechanisms} (peer effects, information cascades, technology spillovers)
\item \textbf{LLM-based modelling tools} (text mining, belief extraction, AI-generated behavioral priors)
\item \textbf{Policy analysis} (Finland’s digitalization strategy, OECD AI metrics, regulatory contexts)
\end{itemize}

The overarching goal is to build an interdisciplinary understanding of how AI technologies diffuse, what accelerates or slows adoption, and how policymakers and organizations can design interventions that support inclusive, trusted, and productive digital transformation.

\section{Research Questions}

The research plan is organized around the following guiding questions:

\subsection*{RQ1: Household Diffusion}

How do individuals and households form beliefs about AI technologies, and how do peer networks shape adoption decisions? In particular, how do online information cascades (e.g., social media and news exposure) versus offline ties (family, workplace, and local communities) differentially affect perceived benefits (productivity), perceived costs (privacy and security), and willingness to adopt?

\noindent \textbf{Method hook:} Survey microdata combined with exposure measures from local and online networks; placebo and alternative exposure definitions for robustness.

\subsection*{RQ2: Firm Digitalization}

How do small and mid-sized Nordic firms decide whether to integrate AI tools such as LLMs, workflow automation, or predictive analytics? Which firm-to-firm network channels (supply chains, local clusters, managerial networks) facilitate spillovers, and what are the key adoption frictions (skills constraints, uncertainty about returns, compliance costs, and organizational change)?

\noindent \textbf{Method hook:} Firm adoption microdata linked to regional/sectoral networks; heterogeneous-treatment and spillover-robust estimators under interference.

\subsection*{RQ3: Network Structure and Inequality}

Do differences in network connectivity (geographical, occupational, or institutional) create persistent gaps in access to or use of AI technologies? Which inequalities are most salient in Nordic contexts—income and education gradients, rural--urban divides, sectoral exposure, or demographic gaps—and can network-targeted policies reduce these disparities?

\noindent \textbf{Method hook:} Inequality decomposition and distributional treatment effects; policy thresholds (where available) exploited for quasi-experimental variation.

\subsection*{RQ4: Policy and Strategy}

How can governments and organizations design interventions that accelerate safe, trustworthy, and inclusive AI adoption? Which interventions are most effective—information and communication strategies, training and re-skilling programs, targeted subsidies, or network-based seeding strategies—in improving measurable outcomes such as adoption rates, productivity proxies, and inequality metrics?

\noindent \textbf{Method hook:} Counterfactual policy simulations and network-targeted seeding experiments in ABM, calibrated to empirically estimated adoption elasticities.

The project combines microdata, network analysis, computational modelling, and policy evaluation to address these questions.

Where feasible, the project will also conduct targeted comparative analyses across Nordic countries using harmonized OECD/EU indicators and public textual corpora to assess how institutional and policy differences moderate diffusion dynamics.

\section{Contribution to the University of Turku}

The project aligns with several research strengths at the University of Turku and is designed to generate concrete institutional benefits, including interdisciplinary collaboration, research training, and contributions to teaching and outreach:

\subsection{Economics}

\begin{itemize}
\item Behavioral microeconomics
\item Applied econometrics
\item Digital transformation
\item Innovation diffusion and productivity
\end{itemize}

\subsection{School of Management}

\begin{itemize}
\item Strategic management and organizational change
\item Technology adoption and innovation systems
\item Data-driven decision-making
\item Human–AI interaction in organizations
\end{itemize}

The project integrates these domains by combining rigorous economic modeling with organizational perspectives, producing insights of direct relevance to Finnish firms and policymakers.

The project will actively foster collaboration through joint seminars and co-authored outputs across Economics and the School of Management, supervision of MSc/PhD student projects (e.g., text-based measures of AI concerns or firm adoption case studies), and the development of teaching materials and case studies related to Finnish and Nordic AI adoption. A realistic target is 1--2 interdisciplinary co-authored papers and 2--3 supervised MSc theses over the project period, subject to departmental fit and student availability.

\section{Theoretical Framework}

The research plan extends classical diffusion and adoption theories by integrating the following topics.

\subsection{Network Economics}
Diffusion modeled through heterogeneous social and organizational structures where individual behavior depends on local neighborhood exposure to digital technologies, building on core network economics and peer effects frameworks (e.g., Jackson, 2010; Manski, 1993).

\subsection{Computational Social Science}
Agent-based models calibrated with Finnish household and firm data to simulate AI adoption under different policy and market conditions, connecting computational models to empirically grounded diffusion mechanisms.

\subsection{LLM-based Behavioral Modeling}
\begin{itemize}
\item NLP extraction of beliefs and sentiment from public discourse (news, surveys, social media).
\item Embedding-based measures of digital readiness and AI-related concerns.
\item LLM-assisted classification of firm strategies and barriers to digitalization.
\end{itemize}

To mitigate concerns about black-box inference and reproducibility, the project will validate LLM-derived measures using a transparent evaluation protocol: (i) human annotation on a stratified subsample of texts, (ii) comparisons with survey-based measures of beliefs where available, and (iii) stability checks across prompts, models, and time.

\subsection{Causal Inference Under Network Spillovers}
Using modern econometric tools (e.g., exposure mappings, spillover-robust treatment effects) to identify causal pathways in diffusion processes, drawing on recent methodological advances in network interference and diffusion empirics.

\section{Methodology}

\subsection{Data Collection}

The project uses multiple datasets:

\begin{itemize}
\item \textbf{OECD AI policy and digital readiness indicators}
\item \textbf{Finnish household digitalization surveys (Statistics Finland)}
\item \textbf{Firm-level digital technology adoption data} (e.g., Statistics Finland business surveys and harmonized Eurostat/OECD measures where appropriate)
\item \textbf{Textual corpora from Nordic news outlets, parliamentary debates, and organizational documents}
\item \textbf{Social network datasets (anonymized) where available}
\end{itemize}

Because access constraints are a primary feasibility risk, the data strategy is explicitly layered: core analyses will rely on publicly available indicators and textual corpora; restricted microdata (e.g., Statistics Finland research datasets) will be pursued via standard research agreements; and network data will be complemented with proxy networks (e.g., co-location, industry similarity, supply-chain linkages) when direct social network data are unavailable.

\subsection{Machine Learning and NLP}

LLMs will be used to:

\begin{itemize}
\item Extract beliefs, expectations, and concerns about AI from Finnish texts
\item Identify diffusion narratives and framing
\item Construct time-varying “AI sentiment indices”
\item Classify firm digitalization strategies
\item Generate behavioral priors for simulation models
\end{itemize}

All NLP pipelines will be implemented with reproducible workflows (versioned corpora, model cards, and documented prompts). Sensitivity analyses will assess the impact of alternative tokenization, model families, and annotation schemes on the resulting indices.

\subsection{Network Analysis}

The project models:

\begin{itemize}
\item Household peer networks (geographical, familial, occupational)
\item Firm interconnections (supply chains, regional clusters, innovation networks)
\end{itemize}

Key metrics include centrality, clustering effects, diffusion bottlenecks, and network resilience.

To strengthen causal interpretation under interference, the project will combine exposure mapping approaches with robustness checks such as placebo exposures, alternative network definitions, and (where feasible) instrumental-variable strategies that exploit plausibly exogenous variation in network exposure.

\subsection{Agent-Based Modelling}

Custom agent-based models will simulate:

\begin{itemize}
\item Household adoption under social influence
\item Firm adoption under uncertainty
\item Policy interventions (subsidies, public communication, education campaigns)
\end{itemize}

Parameters will be calibrated with empirical data and validated using econometric benchmarks, including matching key moments such as adoption rates by firm size/sector/region and survey-based belief distributions.

Implementation will use a Python-based stack (e.g., NetworkX for network objects and Mesa-style agent-based simulation patterns) to facilitate transparency, unit testing, and integration with empirical estimation.

\subsection{Integration Across Methods}

The empirical and computational components are designed to be mutually informative: LLM-extracted measures of beliefs and perceived risks will inform time-varying adoption propensities and heterogeneity in the agent-based models, while econometric estimates of peer effects and spillovers will provide discipline for simulation parameters and serve as validation targets for counterfactual scenarios.

Where feasible and compliant with data agreements, the project will release reproducible code, documented prompts, and non-sensitive derived measures to support transparency and reuse.

\subsection{Policy Evaluation}

The project evaluates Finland’s AI strategy through:

\begin{itemize}
\item Counterfactual simulations
\item Diffusion inequality analysis
\item Network-targeted intervention modelling
\end{itemize}

\section{Secondment: Lucidminds AI (Netherlands)}

Lucidminds AI specializes in agent-based modelling, interactive simulations, and AI-driven narrative environments. The secondment will focus on:

\begin{itemize}
\item Developing explainable agent-based models for AI diffusion
\item Implementing scalable simulation infrastructure
\item Co-designing tools for visualizing network diffusion and policy scenarios
\item Translating academic models into industry-ready prototypes
\end{itemize}

This secondment ensures strong \textbf{intersectoral value}, bridging academic research with technological implementation and stakeholder communication.

\section{Expected Scientific Impact}

\subsection{Advancing Diffusion Theory}

The project provides one of the first frameworks integrating LLM-based behavioral extraction with network diffusion models.

\subsection{Methodological Innovation}

\begin{itemize}
\item Integration of LLM-based belief modelling with economic simulations
\item Novel network-based causal identification strategies
\item Policy evaluation tools with real-world applicability
\end{itemize}

\subsection{High-impact publications}

Target outlets will be tiered to balance ambition and feasibility. Flagship submissions will target top general-interest and field journals, complemented by policy- and interdisciplinary-oriented outlets:

\begin{itemize}
\item \emph{Journal of Economic Behavior \& Organization}
\item \emph{Research Policy}
\item \emph{Management Science}
\item \emph{Journal of Urban Economics}
\item \emph{Journal of the European Economic Association}
\item \emph{Economics of Innovation and New Technology}
\end{itemize}

Additional dissemination-focused outputs (e.g., policy journals and interdisciplinary venues) will support fast translation of results to stakeholders.

\section{Societal Impact and Alignment with Finland’s Strategy}

Finland aims to be “a global leader in human-centric, trusted AI.” The proposed research supports this strategy by providing:

\begin{itemize}
\item Evidence on household barriers to AI adoption
\item Insights into digital skills inequality
\item Tools for designing targeted communication strategies
\item Models predicting diffusion of digital public services
\item Support for OECD and EU-wide digitalization policy monitoring
\end{itemize}

Stakeholders include:

\begin{itemize}
\item Ministry of Economic Affairs and Employment
\item Ministry of Education and Culture
\item OECD Observatory of Public Sector Innovation
\item Nordic Council digitalization policy units
\end{itemize}

Planned outputs include policy briefs, stakeholder workshops, and (where feasible) an interactive dashboard that allows policymakers to explore counterfactual diffusion scenarios and distributional impacts under alternative interventions.

The dashboard will be designed around practical policy levers (e.g., user-adjustable subsidy levels, training intensity, and network targeting rules) and will be paired with a dissemination plan targeting at least 5 relevant policy units and stakeholder groups.

\section{Feasibility and Work Plan}

\subsection{Year 1}

\begin{itemize}
\item Construct datasets
\item Extract beliefs/sentiment using NLP
\item Build initial network maps
\end{itemize}

Year 1 will also include a feasibility pilot using fully public data (text corpora and OECD indicators) to de-risk later stages that depend on restricted microdata.

\subsection{Year 2}

\begin{itemize}
\item Develop and calibrate agent-based diffusion models
\item Conduct secondment at Lucidminds AI
\item Produce first working papers
\end{itemize}

To manage overlap risk between model development and the secondment, core simulation infrastructure will be prototyped in late Year 1 and refined during Year 2.

\subsection{Year 3}

\begin{itemize}
\item Policy simulations and impact evaluation
\item Publish academic articles
\item Disseminate results through TIAS workshops and Nordic policy forums
\end{itemize}

Dissemination in Year 3 will target at least 2--3 presentations/workshops (academic and policy-facing) to ensure uptake by relevant ministries and Nordic policy networks.

\section{Training Needs and Career Development}

The project will strengthen expertise in:

\begin{itemize}
\item Advanced NLP and LLM pipelines
\item Network econometrics
\item Diffusion theory and innovation policy
\item Intersectoral collaboration and stakeholder communication
\end{itemize}

Training will be supported through targeted mentorship and collaboration with UTU faculty in applied econometrics and digital transformation research, complemented by hands-on implementation experience during the Lucidminds secondment.

This fellowship is an ideal steppingstone toward a future academic career focusing on digital economics, AI policy, and computational methods.

\section{Ethical Considerations}

\begin{itemize}
\item No analysis of personal identifiable data
\item All text data anonymized or publicly available
\item LLM applications audited for bias and reproducibility (e.g., documented prompts, subgroup performance checks for text classifiers, and periodic human-in-the-loop review)
\item Compliance with EU AI Act principles
\item Compliance with GDPR and relevant data governance requirements for restricted-access microdata and textual corpora
\item Full adherence to the University of Turku’s ethics and equality policies
\end{itemize}

\nocite{*}

\bibliography{ref}
% \printbibliography

\end{document}
